\begin{scanpage}{0190}{172}
direction toward the source, it may not coincide with the polarization of the wave radiated by the source. Such a phenomenon is called reflected wave depolarization.

It is easy to establish the reason for depolarization, if one investigates the fringing field as the sum of the fields radiated by the uniform and nonuniform parts of the current. According to \S~26, the uniform part of the current radiates the following field in the direction towards the source ($\vartheta=\pi-\gamma$, $\varphi=-\pi/2$)
\[
\left.\begin{aligned}E_x=-H_\vartheta&=\frac{iaE_{0x}}2\frac{e^{ikR}}R\overline\Sigma^0,\\
E_\vartheta=H_\varphi&=\frac{iaH_{0x}}2\frac{e^{ikR}}R\Sigma^0.
\end{aligned}\right\}\tag{27.04}
\]
The functions $\overline\Sigma^0$ and $\Sigma^0$ satisfy the condition $\Sigma^0=-\overline\Sigma^0$, and are described by Equation (26.15). From Equation (27.04), let us immediately obtain the equality
\[
\frac{E_x}{E_\vartheta}=\operatorname{tg}\alpha,\tag{27.05}
\]
which means that in the physical optics approach the reflected wave does not experience depolarization. Consequently, the reflected wave depolarization is caused \underline{only by the nonuniform part of the current} or, in other words, by the surface distortion.

Let us derive an equation for the magnitude of angle $\delta$. This is the angle by which the electric field vector of the reflected wave is turned in respect to the electric vector of the wave radiated by the source. For this purpose, let us represent the functions $\overline\Sigma_{1(2)}$ and $\Sigma_{1(2)}$ in the form
\[
\left.\begin{aligned}\overline\Sigma_{1(2)}&=\overline\Sigma^0_{1(2)}+\overline\Sigma^1_{1(2)},\\
\Sigma_{1(2)}&=\Sigma^0_{1(2)}+\Sigma^1_{1(2)},
\end{aligned}\right\}\tag{27.06}
\]
where the terms $\overline\Sigma^0_{1(2)}$ and $\Sigma^0_{1(2)}$ correspond to the field radiated by the uniform part of the current, and the terms $\overline\Sigma^1_{1(2)}$ and $\Sigma^1_{1(2)}$
\end{scanpage}
